\section{Compactness And Exact-Object Capture}
\label{sec:compactness-exact-object-branch}

The compactness branch addresses a different part of the positive-forward proof
program. Energy, enstrophy, and scale estimates try to produce control. The
compactness branch asks what kind of object that control produces near an alleged
terminal time. A finite-time breakdown claim needs more than a growing norm. It
needs a terminal object that still belongs to the original Navier--Stokes
problem.

The positive-forward aim is exact-object capture. Start with a smooth solution
on its maximal interval \([0,T_*)\), and suppose a terminal scenario is being
formed as \(t\uparrow T_*\). One selects a sequence \(t_n\uparrow T_*\) and a
family of localized or spectral objects derived from \(u(t_n)\). Compactness is
then used to extract a limit. The useful limit must remain attached to the same
datum, the same equation, the same pressure law, and the same preterminal
solution.

The standard compactness pattern has two parts. Low frequencies become compact
once enough time regularity and spatial control are available. High frequencies
are made small by a tail estimate. Schematically, write
\[
  u=P_{\leq N}u+P_{>N}u.
\]
The low-frequency term lives in a finite-dimensional or locally compact space
after equicontinuity is obtained. The high-frequency term is controlled by a
spectral tail quantity such as
\[
  \sum_{j>N}2^{2sj}\|P_j u(t)\|_{L^2}^2.
\]
The compactness plan is to choose \(N\) large, control the tail uniformly, pass to
a limit in the low modes, and reconstruct the terminal object.

The branch stalls at dependency order. Compactness consumes exactly the controls
that the earlier branches were trying to prove. A low/high split can extract an
object only after the tail, source, and participation bills have already been
paid at the strength required by the topology. Thus compactness is a receiver of
control rather than a free source of control.

This creates the exact-object residue. A terminal limit can be too weak, too
detached, too dependent on the chosen sequence, or too poorly connected to the
pressure-viscosity law. Each defect matters because a breakdown witness has to be
a same-solution witness. A detached compactness artifact cannot carry the full
burden of finite-time breakdown for the original Navier--Stokes solution.

The compactness branch therefore supplies one of the central motivations for the
CM-exit method. A terminal object has to be tested before it is allowed to serve
as a witness. The test is not merely whether some limiting distribution exists.
The test asks whether the object carries the requirements for membership in
the working Navier--Stokes class.

Three failure modes dominate this branch.

\begin{enumerate}[leftmargin=*]
\item \textbf{Carrier failure.} The selected object lacks a positive material or
      packet carrier at the terminal scale. This is a Pack-side obstruction.
\item \textbf{Participation failure.} The selected object fails to carry the same
      pressure-viscosity participation law. This is a Part-side obstruction.
\item \textbf{Readout failure.} The selected object has no positive-scale field
      readout strong enough to support continuation. This is a Field-side
      obstruction.
\end{enumerate}

The positive-forward compactness branch tries to produce a terminal object that
survives all three tests. The CM-exit method reverses the burden at the terminal
claim. If a finite-time breakdown claim points to a compactness object, that
object must pass the same-solution membership requirements. The first failed requirement
turns the object into a class-exit witness rather than an in-class nonsmooth
continuation.

This branch also explains why the paper must track surface fidelity. A compact
limit on the torus, a local blowup profile, and a whole-space export object are
not interchangeable. Each lives on a specific surface with specific compactness
and pressure rules. The paper treats periodic packet capture,
local profile capture, and whole-space export as related branch concerns with
separate admission tests.
